\documentclass[journal]{IEEEtran}

\usepackage{amsmath,amssymb,amsfonts}
\usepackage{algorithmic}
\usepackage{graphicx}
\usepackage{textcomp}
\usepackage{xcolor}
\usepackage{booktabs}
\usepackage{multirow}
\usepackage{url}
\usepackage{hyperref}
\usepackage{microtype}
\usepackage{cite}

\hypersetup{
    colorlinks=true,
    linkcolor=blue,
    citecolor=blue,
    urlcolor=blue
}

\begin{document}

\title{Stress-Testing Longitudinal Synthetic EHR and ICU Waveform Generation: An Adversarial Benchmarking Framework for Privacy, Utility, and Algorithmic Fairness}

\author{Aaqib~Rashid~Mir,~\IEEEmembership{Member,~IEEE}%
\thanks{A. R. Mir is with the Department of Artificial Intelligence \& Data Science, Guru Nanak Institutions Technical Campus (Autonomous), Hyderabad, Telangana 501506, India, and also with The Coding Science --- School of Technology \& AI Innovations (e-mail: mtechbro94@gmail.com). ORCID: 0009-0002-3921-654X.}%
}

\markboth{IEEE Transactions on Biomedical Engineering / IEEE JBHI,~Vol.~XX, No.~XX,~2026}%
{Mir: Stress-Testing Longitudinal Synthetic EHR and ICU Waveform Generation}

\maketitle

\begin{abstract}
Deep generative architectures---spanning continuous-time diffusion models, generative adversarial networks (GANs), and longitudinal autoregressive transformers---are widely championed as a privacy-preserving paradigm to liberate restricted Intensive Care Unit (ICU) Electronic Health Records (EHR) for secondary scientific investigation. The operational presumption underpinning this movement is that synthetic trajectories decouple clinical predictive utility from re-identification liability. However, prevailing evaluation regimes remain fundamentally fragmented: they almost universally rely on low-dimensional static marginal distributions and aggregate similarity metrics, which systematically fail to detect non-physiological temporal dynamics or empirical privacy leakage under directed adversarial scrutiny. In this paper, we establish an end-to-end, empirical stress-testing benchmark designed specifically for longitudinal synthetic EHR and physiological waveforms along three interconnected clinical trust axes: \textbf{Privacy}, \textbf{Utility}, and \textbf{Fairness}. Our framework introduces: (1) a four-way partitioned shadow-model Membership Inference Attack (MIA) and sensitive attribute reconstruction pipeline tailored to continuous multivariate time-series; (2) dual-direction clinical utility benchmarks implementing Train-on-Synthetic-Test-on-Real (TSTR) and Train-on-Real-Test-on-Synthetic (TRTS) pipelines targeting early patient deterioration; (3) five formal temporal realism probes that quantify sequence irregularity, non-uniform measurement cadence, first-order Markovian state transition matrix drift over discretized physiological states, within-stay autocorrelation decay error, and cumulative multiday trend drift; and (4) demographic subgroup parity audits quantifying disparate privacy vulnerability and utility gaps across marginalized patient sub-cohorts. We benchmark canonical generative paradigms (Gaussian Copula, 1D-CNN WGAN-GP, TimeGAN, and TabDDPM-Time) on a curated cohort of 42,110 adult ICU admissions from the MIMIC-IV database (v2.2). While unconstrained models demonstrated high nominal downstream utility (TabDDPM achieving TSTR AUROC $0.814$ [$95\%$ CI: $0.802$--$0.826$] compared to the real-data baseline of $0.842$ [$0.831$--$0.853$]), this was accompanied by significant empirical privacy leakage (MIA AUROC $0.618$, adversary advantage $0.236$). Recurrent adversarial networks exhibited catastrophic temporal drift: WGAN-GP and TimeGAN displayed transition matrix divergence ($D_{\text{trans}} = 0.224$ and $0.148$), generating physiologically impossible hourly jumps between extreme bradycardia ($<60$~bpm) and malignant tachycardia ($>120$~bpm) at rates over 28-fold higher than observed in clinical reality ($P_{\text{syn}} = 0.084$ vs. $P_{\text{real}} = 0.003$). Crucially, demographic auditing uncovered severe asymmetric vulnerabilities: minority racial cohorts and elderly patients ($\ge 65$ years) suffered up to an $18.2\%$ increase in membership re-identification vulnerability and an $8.4\%$ drop in predictive transfer utility relative to the majority population. Our findings demonstrate that static distributional alignment masks profound temporal instability and disproportionate privacy risks for vulnerable patient populations.
\end{abstract}

\begin{IEEEkeywords}
Synthetic Healthcare Data, Electronic Health Records, Differential Privacy, Membership Inference Attack, Clinical Utility, Algorithmic Fairness, MIMIC-IV, Time-Series Modeling.
\end{IEEEkeywords}

\section{Introduction}
\IEEEPARstart{T}{he} secondary analysis of granular Electronic Health Records (EHR) and continuous physiological waveform streams acquired in intensive care units (ICUs) has propelled modern clinical informatics, enabling major breakthroughs in algorithmic early-warning systems, dynamic acuity estimation, and personalized hemodynamic stabilization \cite{johnson2023mimic, rajkomar2018scalable}. Despite these methodological strides, medical artificial intelligence research remains encumbered by an institutional data-access paradox. Regulatory protections---such as the Health Insurance Portability and Accountability Act (HIPAA) in the United States and the General Data Protection Regulation (GDPR) in the European Union---properly restrict the unfettered dissemination of clinical trajectories to prevent the re-identification of vulnerable patients \cite{price2019privacy}. Consequently, the vast majority of biomedical predictive architectures are developed, tuned, and validated on isolated, single-institution data silos, hampering cross-institutional reproducibility and exacerbating algorithmic bias.

To circumvent the logistical, legal, and institutional liabilities of raw data sharing, deep generative modeling has emerged as a prominent paradigm \cite{choi2017generating, yoon2019timegan}. Generative Adversarial Networks (GANs), score-based continuous diffusion models, and temporal transformer decoders are trained to approximate the high-dimensional probability manifold of real patient records. The synthetic records sampled from these distributions are frequently posited to be intrinsically anonymous---operating under the assumption that because a generated trajectory does not correspond to any singular physical individual, it is immune to regulatory embargoes and free from privacy concerns \cite{hernandez2022synthetic}.

\subsection{The Failure Modes of Contemporary Validation Regimes}
Despite substantial investment in synthetic health data generation, the path from experimental algorithms to trusted scientific infrastructure is impeded by an acute methodological failure: the lack of rigorous, multidimensional, and domain-grounded evaluation standards. The prevailing literature overwhelmingly assesses synthetic EHR quality through narrow statistical lenses, predominantly comparing static one-dimensional marginal distribution histograms, low-order pairwise correlation matrices, or aggregate distance metrics such as Maximum Mean Discrepancy (MMD) \cite{hernandez2022synthetic, alaa2022generative}.

These static statistical metrics obscure three structural clinical failure modes:
\begin{enumerate}
    \item \textbf{The Temporal Plausibility Illusion:} Medical time-series data do not consist of independent and identically distributed (i.i.d.) observations; rather, they reflect continuous homeostatic feedback loops governed by physiological laws. A synthetic vital sign trajectory can match empirical population-level means, standard deviations, and inter-quartile ranges with exceptional precision while generating sequences that are biologically impossible. For example, a generative model may produce heart rate trajectories that swing from severe bradycardia ($<45$~beats/min) to supraventricular tachycardia ($>150$~beats/min) and back across consecutive hourly observations without accompanying changes in systemic blood pressure or medication administration. Such temporal discoherence destroys clinical credibility and can cause downstream predictive models to learn spurious artifacts.
    \item \textbf{The Fragility of Empirical Privacy:} While theoretical privacy mechanisms such as $(\epsilon, \delta)$-Differential Privacy (DP) provide provable upper bounds on worst-case disclosure risk \cite{dwork2014algorithmic}, practical implementations of Differentially Private Stochastic Gradient Descent (DP-SGD) in continuous, multi-channel longitudinal domains frequently suffer catastrophic utility degradation \cite{abadi2016deep}. In response, practitioners routinely deploy unconstrained generative architectures under the unverified assumption that neural generalization alone will prevent verbatim memorization. In reality, deep generative models trained on complex, high-dimensional manifolds are notoriously prone to memorizing isolated outlier trajectories---precisely the critically ill, complex patient presentations that characterize high-stakes ICU admissions \cite{carlini2019secret, stadler2022synthetic}.
    \item \textbf{Demographic Disparity Amplification (Fairness Erosion):} Real-world clinical datasets exhibit deep structural demographic skews, with substantial underrepresentation of specific racial, ethnic, and elderly subpopulations \cite{obermeyer2019dissecting}. Standard generative optimization objectives (such as minimizing the Wasserstein distance or evidence lower bound) inherently prioritize global manifold approximation, optimizing generative fidelity for dominant demographic cohorts at the expense of distribution tails. This results in severe ``fairness gaps,'' wherein synthetic records for minority cohorts exhibit degraded clinical predictive utility alongside heightened vulnerability to adversarial membership re-identification \cite{bagdasaryan2019differential}.
\end{enumerate}

\subsection{Contributions of This Work}
To address these systemic vulnerabilities, we introduce the \textbf{Synthetic EHR Benchmark}, a standardized empirical framework designed to evaluate longitudinal synthetic clinical records and ICU waveforms along three essential axes: Privacy, Utility, and Fairness.

The primary methodological and empirical contributions of this study are:
\begin{itemize}
    \item \textbf{A Standardized, Multi-Vector Threat Model:} We formalize shadow-model membership inference attacks (MIA) and sensitive attribute inference attacks specifically parameterized for multi-channel temporal trajectories, evaluating adversary advantage and precision at strict false-positive thresholds.
    \item \textbf{Dual-Transfer Clinical Utility Pipelines:} We operationalize Train-on-Synthetic-Test-on-Real (TSTR) and Train-on-Real-Test-on-Synthetic (TRTS) benchmarking protocols across canonical ICU risk-stratification endpoints: 24-hour mortality and acute hemodynamic decompensation.
    \item \textbf{Five Physiologically Grounded Temporal Realism Probes:} We develop a suite of temporal probes designed to quantify trajectory-length divergence, measurement interval regularity, first-order Markovian state transition matrix drift over discretized physiological bins, within-stay autocorrelation decay error, and cumulative multiday trend drift.
    \item \textbf{Subgroup Fairness Stress Testing:} We formulate fairness audit protocols that quantify disparities in downstream utility and empirical re-identification risk across intersecting demographic partitions (biological sex, age brackets, and racial/ethnic groups).
    \item \textbf{Comprehensive Benchmarking on MIMIC-IV:} We execute an empirical investigation evaluating four canonical generative paradigms (Gaussian Copula, WGAN-GP, TimeGAN, and TabDDPM-Time) over $42,110$ admissions from the MIMIC-IV database, documenting performance trade-offs, temporal failure modes, and equity disparities. All code and evaluation routines are released as open-source software.
\end{itemize}

\section{Related Work}
\subsection{Deep Generative Models for Clinical Time-Series}
The synthesis of electronic health records has evolved rapidly from early parametric copula formulations and Bayesian networks to deep neural architectures \cite{choi2017generating, tucker2020generating}. MedGAN \cite{choi2017generating} demonstrated the feasibility of generating discrete diagnostic billing codes by coupling a multi-layer perceptron GAN with a pre-trained autoencoder. Subsequent frameworks addressed intra-record correlation; CorGAN \cite{torfi2020corgan} introduced convolutional autoencoders to capture cross-code associations, while Choi et al. \cite{choi2016mime} explored recurrent neural networks (RNNs) for multi-visit sequences.

For continuous physiological signals and dense ICU vitals, TimeGAN \cite{yoon2019timegan} unified autoregressive sequence modeling with adversarial objectives, enforcing temporal dynamics through an explicit stepwise latent transition loss. More recently, score-based continuous-time diffusion models---exemplified by TabDDPM \cite{kotelnikov2023tabddpm} and its temporal variants---have established superior sample quality in tabular and sequential domains by modeling synthesis as a reverse stochastic differential equation. Concurrently, masked autoencoding transformers have been adapted for multi-channel clinical trajectories \cite{hur2022unifying}. However, the biomedical generative literature remains overwhelmingly focused on model architecture and generative fidelity, while formal adversarial stress-testing protocols tailored to longitudinal ICU dynamics remain largely unstandardized.

\subsection{Privacy Leakage and Adversarial Attacks in Medical AI}
The concept that synthetic data is intrinsically private has been repeatedly debunked across machine learning literature. Shokri et al. \cite{shokri2017membership} introduced shadow-model membership inference attacks (MIA), proving that machine learning classifiers leak the presence of individual training samples through overconfident prediction vectors. Chen et al. \cite{chen2020gan} and Hayes et al. \cite{hayes2019logan} extended MIA methodologies to generative models, establishing that generative networks leak membership when the generated data distribution aligns too tightly with localized training manifolds.

In healthcare specifically, Stadler et al. \cite{stadler2021privacy, stadler2022synthetic} conducted an extensive audit of synthetic tabular health records, demonstrating that synthetic datasets offering competitive downstream utility frequently exhibit elevated membership re-identification risks. Longitudinal continuous waveforms amplify this vulnerability by providing high-dimensional, unique temporal signatures \cite{esteban2017real}. Despite the theoretical protections offered by differential privacy \cite{dwork2014algorithmic}, empirical studies by Bagdasaryan et al. \cite{bagdasaryan2019differential} have shown that applying DP-SGD in complex settings often incurs disproportionately severe utility degradation on minority demographic cohorts.

\subsection{Evaluating Clinical Utility and Temporal Coherence}
The conceptual foundation for empirical utility assessment in synthetic data was formalized by Esteban et al. \cite{esteban2017real} through the Train-on-Synthetic, Test-on-Real (TSTR) framework. Under this paradigm, a supervised model trained purely on synthetic trajectories is evaluated on real holdout test records. A minimal discrepancy between TSTR and the real-data baseline (Train-on-Real, Test-on-Real; TRTR) serves as an empirical indicator of utility preservation.

However, an emerging body of critical literature highlights that TSTR is insensitive to subtle dynamical and physiological implausibilities \cite{hernandez2022synthetic, theis2016note}. A clinical risk model can achieve high discriminative performance on static averages or linear proxies while remaining oblivious to impossible physiological oscillations within the synthetic training set. To date, literature lacks a cohesive framework that couples TSTR utility evaluation with rigorous, physics- and physiology-grounded temporal realism probes.

\section{Problem Formulation and Threat Models}

\subsection{Mathematical Notation and Data Representation}
Let $\mathcal{D}_{\text{real}} = \{(\mathbf{X}_i, \mathbf{y}_i, \mathbf{s}_i)\}_{i=1}^{N_{\text{real}}}$ denote a cohort of real patient admissions drawn from an underlying population distribution $\mathcal{P}_{\text{clinical}}$. Each patient record $i$ is characterized by:
\begin{itemize}
    \item A continuous, multivariate longitudinal trajectory $\mathbf{X}_i \in \mathbb{R}^{T_i \times D}$, where $T_i$ represents the number of sequential temporal observations and $D$ represents the dimension of monitored physiological channels.
    \item A vector of binary clinical outcome labels $\mathbf{y}_i \in \{0, 1\}^K$, denoting critical downstream events such as in-hospital mortality or acute clinical decompensation.
    \item A vector of sensitive demographic attributes $\mathbf{s}_i \in \mathcal{S}$, encoding categorical variables such as biological sex, age bracket, and self-reported race/ethnicity.
\end{itemize}

A parameterized generative model $G_\theta: \mathcal{Z} \to \mathcal{X}$, conditioned on parameters $\theta$ and driven by a standard latent prior $\mathbf{z} \sim \mathcal{N}(\mathbf{0}, \mathbf{I})$, produces a synthetic clinical cohort:
\begin{equation}
\mathcal{D}_{\text{syn}} = \{(\tilde{\mathbf{X}}_j, \tilde{\mathbf{y}}_j, \tilde{\mathbf{s}}_j)\}_{j=1}^{N_{\text{syn}}} \sim \mathcal{P}_{G_\theta}
\end{equation}

Our benchmarking objective is to formally evaluate the empirical risk surface of $G_\theta$ across the tripartite axes of privacy, utility, and fairness.

\subsection{Formal Threat Model: Membership Inference Attacks (MIA)}
We define an adversarial game under the standard \textit{shadow-model framework} \cite{shokri2017membership, stadler2022synthetic}. The adversary $\mathcal{A}$ seeks to determine whether a specific patient trajectory $\mathbf{x}^* = (\mathbf{X}^*, \mathbf{y}^*, \mathbf{s}^*)$ was an active participant in the training partition $\mathcal{D}_{\text{train}} \subset \mathcal{D}_{\text{real}}$ utilized to fit the target generator $G_\theta$.

We adopt a black-box, worst-case threat model adhering to the following axioms:
\begin{enumerate}
    \item \textbf{Access to Synthetic Data:} The adversary has full inspection access to the released synthetic dataset $\mathcal{D}_{\text{syn}}$.
    \item \textbf{Background Knowledge:} The adversary possesses an auxiliary dataset $\mathcal{D}_{\text{aux}}$ drawn from the same clinical catchment population $\mathcal{P}_{\text{clinical}}$, such that $\mathcal{D}_{\text{aux}} \cap \mathcal{D}_{\text{train}} = \emptyset$.
    \item \textbf{Shadow Synthesis:} The adversary partitions $\mathcal{D}_{\text{aux}}$ into disjoint subsets $\mathcal{D}_{\text{shadow}}^{\text{in}}$ and $\mathcal{D}_{\text{shadow}}^{\text{out}}$, trains one or more shadow generators $G_\phi$ on $\mathcal{D}_{\text{shadow}}^{\text{in}}$, and generates shadow synthetic cohorts $\mathcal{D}_{\text{shadow}}^{\text{syn}}$.
\end{enumerate}

To operationalize the attack across multi-channel temporal trajectories, the adversary computes an identification feature vector $\mathbf{v}(\mathbf{x}^*, \mathcal{D}_{\text{syn}}) \in \mathbb{R}^M$. For a candidate record $\mathbf{x}^*$, this vector encodes local density and proximity relative to the synthetic manifold:
\begin{enumerate}
    \item \textbf{Minimum Trajectory Euclidean Distance:}
    \begin{equation}
    d_{\min}(\mathbf{x}^*, \mathcal{D}_{\text{syn}}) = \min_{\tilde{\mathbf{x}} \in \mathcal{D}_{\text{syn}}} \|\mathbf{x}^* - \tilde{\mathbf{x}}\|_F
    \end{equation}
    where $\|\cdot\|_F$ denotes the Frobenius norm across time and channels.
    \item \textbf{$k$-Nearest Neighbor Proximity:}
    \begin{equation}
    d_k(\mathbf{x}^*, \mathcal{D}_{\text{syn}}) = \frac{1}{k} \sum_{\tilde{\mathbf{x}} \in \mathcal{N}_k(\mathbf{x}^*; \mathcal{D}_{\text{syn}})} \|\mathbf{x}^* - \tilde{\mathbf{x}}\|_F
    \end{equation}
    \item \textbf{Dynamic Time Warping (DTW) Phase-Aligned Distance:}
    \begin{equation}
    \text{DTW}(\mathbf{x}^*, \tilde{\mathbf{x}}) = \min_{\pi} \sqrt{\sum_{(t_1, t_2) \in \pi} \|x^*_{t_1, :} - \tilde{x}_{t_2, :}\|_2^2}
    \end{equation}
    \item \textbf{Kernel Density Manifold Estimation:}
    \begin{equation}
    \hat{\rho}(\mathbf{x}^*; \mathcal{D}_{\text{syn}}, \sigma) = \frac{1}{|\mathcal{D}_{\text{syn}}|} \sum_{\tilde{\mathbf{x}} \in \mathcal{D}_{\text{syn}}} \exp\left(-\frac{\|\mathbf{x}^* - \tilde{\mathbf{x}}\|_F^2}{2\sigma^2}\right)
    \end{equation}
\end{enumerate}

The attack model $\mathcal{A}: \mathbf{v}(\mathbf{x}, \mathcal{D}_{\text{syn}}) \to [0, 1]$ is trained as a gradient-boosted decision ensemble on labeled shadow vectors. Performance is evaluated via:
\begin{itemize}
    \item \textbf{Area Under the Receiver Operating Characteristic Curve ($\text{AUROC}_{\text{MIA}}$):}
    \begin{equation}
    \text{AUROC}_{\text{MIA}} = \int_0^1 \text{TPR}(\text{FPR}) \, d(\text{FPR})
    \end{equation}
    \item \textbf{Adversary Advantage ($\text{Adv}$):}
    \begin{equation}
    \text{Adv} = \max_{\tau} |\text{TPR}(\tau) - \text{FPR}(\tau)| = 2 \times |\text{AUROC}_{\text{MIA}} - 0.50|
    \end{equation}
    \item \textbf{Precision at Strict False Positive Rate ($\text{Precision@1\%FPR}$)}.
\end{itemize}

\subsection{Formal Threat Model: Attribute Inference Attacks (AIA)}
In the Attribute Inference Attack (AIA) formulation, the adversary possesses a partially observed patient trajectory $\mathbf{X}^*_{\text{obs}} \in \mathbb{R}^{T \times D_{\text{obs}}}$ and seeks to predict a sensitive, withheld clinical or demographic attribute $s^* \in \mathcal{S}$.

The adversary leverages $\mathcal{D}_{\text{syn}}$ to fit an attribute decoder $f_{\text{AIA}}: \mathbb{R}^{T \times D_{\text{obs}}} \to \mathcal{S}$ by solving the empirical risk minimization problem:
\begin{equation}
\hat{\psi} = \arg\min_\psi \frac{1}{|\mathcal{D}_{\text{syn}}|} \sum_{j=1}^{|\mathcal{D}_{\text{syn}}|} \mathcal{L}_{\text{ce}}\left(f_\psi(\tilde{\mathbf{X}}_{j, \text{obs}}), \tilde{s}_j\right)
\end{equation}

The trained decoder is deployed against real patient test records $\mathcal{D}_{\text{real}}^{\text{test}}$. We quantify attribute leakage via the \textbf{Attacker Attribute Lift} ($\Delta_{\text{attr}}$):
\begin{equation}
\Delta_{\text{attr}} = \text{Accuracy}\left(f_{\hat{\psi}}(\mathbf{X}_{\text{real}, \text{obs}}), s_{\text{real}}\right) - \max_{c \in \mathcal{S}} \mathbb{P}(s = c)
\end{equation}

\subsection{Formal Clinical Utility Transfer: TSTR and TRTS}
To verify clinical utility, we deploy the dual-transfer paradigm:
\begin{enumerate}
    \item \textbf{Train on Synthetic, Test on Real (TSTR):}
    A supervised clinical model $f_{\text{TSTR}}$ is optimized over $\mathcal{D}_{\text{syn}}$ and evaluated on held-out real data $\mathcal{D}_{\text{real}}^{\text{test}}$. The empirical TSTR gap relative to the gold-standard baseline (Train on Real, Test on Real; TRTR) is:
    \begin{equation}
    \Delta_{\text{TSTR}} = \mathcal{M}(f_{\text{TRTR}}; \mathcal{D}_{\text{real}}^{\text{test}}) - \mathcal{M}(f_{\text{TSTR}}; \mathcal{D}_{\text{real}}^{\text{test}})
    \end{equation}
    where $\mathcal{M}$ represents a standard discrimination metric (AUROC or AUPRC).
    \item \textbf{Train on Real, Test on Synthetic (TRTS):}
    Conversely, a reference model $f_{\text{TRTS}}$ trained on $\mathcal{D}_{\text{real}}^{\text{train}}$ is evaluated on $\mathcal{D}_{\text{syn}}$. Discrepancies between TRTR and TRTS indicate regions where the synthetic data distribution deviates from authentic physiological response boundaries.
\end{enumerate}

\subsection{Formal Temporal Realism Probes}
To evaluate physiological plausibility, we formulate five temporal probes:

\subsubsection{Probe 1: Sequence Length Distribution Distance ($D_{\text{len}}$)}
Using the Wasserstein-1 metric:
\begin{equation}
D_{\text{len}} = \mathcal{W}_1(P_{\text{real}}, P_{\text{syn}}) = \int_{-\infty}^\infty |F_{\text{real}}(t) - F_{\text{syn}}(t)| \, dt
\end{equation}

\subsubsection{Probe 2: Sampling Interval Discrepancy ($D_{\text{rate}}$)}
Measured via the Jensen-Shannon Divergence ($\text{JSD}$):
\begin{equation}
D_{\text{rate}} = \text{JSD}(P_{\text{real}}(\Delta t) \parallel P_{\text{syn}}(\Delta t))
\end{equation}

\subsubsection{Probe 3: First-Order Markovian State Transition Drift ($D_{\text{trans}}$)}
Discretizing each channel into physiological states $\Omega = \{\text{Hypo}, \text{Normal}, \text{Hyper}\}$, we compute empirical transition probability matrices $\mathbf{P}_d^{\text{real}}, \mathbf{P}_d^{\text{syn}} \in \mathbb{R}^{3 \times 3}$. The drift is the Frobenius norm across all channels:
\begin{equation}
D_{\text{trans}} = \frac{1}{D} \sum_{d=1}^D \|\mathbf{P}_d^{\text{real}} - \mathbf{P}_d^{\text{syn}}\|_F
\end{equation}

\subsubsection{Probe 4: Within-Stay Autocorrelation Decay Error ($D_{\text{autocorr}}$)}
For channel $d$, the autocorrelation at lag $\tau$ is $\rho_d(\tau)$. The aggregate error up to lag $\tau_{\max} = 12$ hours is:
\begin{equation}
D_{\text{autocorr}} = \frac{1}{D \cdot \tau_{\max}} \sum_{d=1}^D \sum_{\tau=1}^{\tau_{\max}} |\rho_d^{\text{real}}(\tau) - \rho_d^{\text{syn}}(\tau)|
\end{equation}

\subsubsection{Probe 5: Multiday Cumulative Physiological Drift ($D_{\text{drift}}$)}
Evaluates directional trajectory drift by calculating the Wasserstein-1 distance between the first-difference velocity distributions $\Delta x_{t, d} = x_{t+1, d} - x_{t, d}$ across consecutive 12-hour epochs.

\subsection{Formal Demographic Fairness Formulations}
Let $\mathcal{S} = \{g_1, g_2, \dots, g_M\}$ represent protected demographic cohorts. We define:
\begin{enumerate}
    \item \textbf{Subgroup Utility Disparity ($\Delta_{\text{fairness}}^{\text{util}}$):}
    \begin{equation}
    \Delta_{\text{fairness}}^{\text{util}} = \max_{g \in \mathcal{S}} \text{AUROC}_{\text{TSTR}}(g) - \min_{g \in \mathcal{S}} \text{AUROC}_{\text{TSTR}}(g)
    \end{equation}
    \item \textbf{Subgroup Privacy Disparity ($\Delta_{\text{fairness}}^{\text{priv}}$):}
    \begin{equation}
    \Delta_{\text{fairness}}^{\text{priv}} = \max_{g \in \mathcal{S}} \text{AUROC}_{\text{MIA}}(g) - \min_{g \in \mathcal{S}} \text{AUROC}_{\text{MIA}}(g)
    \end{equation}
    \item \textbf{Fairness-Privacy Pareto Metric ($\Phi$):}
    \begin{equation}
    \Phi = \frac{\min_{g \in \mathcal{S}} \text{AUROC}_{\text{TSTR}}(g)}{1.0 + \left(\max_{g \in \mathcal{S}} \text{AUROC}_{\text{MIA}}(g) - 0.50\right)}
    \end{equation}
\end{enumerate}

\section{Cohort Construction and Experimental Setup}

\subsection{MIMIC-IV Extraction Protocol}
Empirical evaluation was conducted using the \textbf{MIMIC-IV} clinical database (v2.2) \cite{johnson2023mimic}. We extracted adult ICU encounters according to the following inclusion criteria:
\begin{itemize}
    \item Age $\ge 18$ years at the time of index ICU admission.
    \item Total ICU length of stay $\ge 24$ hours. For patients with multiple ICU admissions, only the initial encounter was included.
    \item Standardized observation window of 24 hours ($T = 24$), sampled at 1-hour intervals.
\end{itemize}

Six continuous vital sign channels were extracted:
\begin{enumerate}
    \item Heart Rate (HR, itemid 220045, beats/min)
    \item Systolic Blood Pressure (SBP, itemid 220050, mmHg)
    \item Diastolic Blood Pressure (DBP, itemid 220051, mmHg)
    \item Respiratory Rate (RR, itemid 220210, breaths/min)
    \item Peripheral Oxygen Saturation ($\text{SpO}_2$, itemid 220277, \%)
    \item Body Temperature (Temp, itemid 223761, $^\circ\text{C}$)
\end{enumerate}

Extreme non-physiological artifacts were removed using consensus clinical ranges. Missing values were imputed via forward-fill followed by population median substitution; a parallel binary missingness mask $\mathbf{M} \in \{0, 1\}^{24 \times 6}$ was appended to preserve sparsity context.

The final cohort comprises \textbf{42,110 adult ICU encounters} ($54.2\%$ Male, $45.8\%$ Female; median age $65.4$ years, IQR: $53.2$--$77.1$; $68.2\%$ White, $11.4\%$ Black/African American, $4.8\%$ Hispanic/Latino, $3.9\%$ Asian, $11.7\%$ Other/Unknown). Overall in-hospital mortality was $10.8\%$ ($n = 4,548$).

\subsection{Benchmark Baseline Architectures}
We evaluated four foundational generative paradigms:
\begin{enumerate}
    \item \textbf{Multivariate Gaussian Copula:} Classical parametric baseline modeling marginal distributions via empirical estimation with a joint Gaussian copula.
    \item \textbf{Continuous WGAN-GP (1D-CNN):} Wasserstein GAN with Gradient Penalty ($\lambda = 10.0$) using 4-layer 1D convolutional residual blocks.
    \item \textbf{TimeGAN \cite{yoon2019timegan}:} Adversarial recurrent framework combining embedding, generator, discriminator, and recovery networks.
    \item \textbf{Score-Based Diffusion (TabDDPM-Time) \cite{kotelnikov2023tabddpm}:} Continuous diffusion framework utilizing a 1D U-Net backbone with residual temporal self-attention.
\end{enumerate}

All deep models were implemented in PyTorch 2.1 and trained on NVIDIA A100 (80GB) GPUs with identical optimization budgets (100 epochs, Adam optimizer, learning rate $2 \times 10^{-4}$, batch size 128). Downstream utility was evaluated using gradient-boosted decision trees (LightGBM) with 5-fold cross-validation. All metrics are reported with 95\% bootstrap confidence intervals from 1,000 stratified resamples.

\begin{table*}[t]
\centering
\caption{Consolidated Benchmark Results Across Privacy, Utility, Temporal Realism, and Fairness on MIMIC-IV. Values in Brackets Indicate 95\% Bootstrap Confidence Intervals. Bold Entries Indicate Superior Synthetic Performance.}
\label{tab:main_results}
\begin{tabular}{lccccccccc}
\toprule
\textbf{Generative Architecture} & \textbf{MIA AUROC} $\downarrow$ & \textbf{Advantage} $\downarrow$ & \textbf{Attr. Acc.} $\downarrow$ & \textbf{TSTR AUROC} $\uparrow$ & \textbf{TRTS AUROC} $\uparrow$ & \textbf{$D_{\text{trans}}$} $\downarrow$ & \textbf{$D_{\text{autocorr}}$} $\downarrow$ & \textbf{$\Delta_{\text{fairness}}^{\text{util}}$} $\downarrow$ & \textbf{Pareto $\Phi$} $\uparrow$ \\
\midrule
\textbf{Real Reference (TRTR)} & $0.500$ \scriptsize{[0.492--0.508]} & $0.000$ & $0.500^*$ & $0.842$ \scriptsize{[0.831--0.853]} & $0.842$ \scriptsize{[0.831--0.853]} & $0.000$ & $0.000$ & $0.038$ \scriptsize{[0.026--0.050]} & $0.842$ \\
\textbf{Gaussian Copula} & \textbf{0.508} \scriptsize{[0.499--0.517]} & \textbf{0.016} & \textbf{0.531} \scriptsize{[0.520--0.542]} & $0.621$ \scriptsize{[0.608--0.634]} & $0.598$ \scriptsize{[0.585--0.611]} & $0.412$ \scriptsize{[0.398--0.426]} & $0.384$ \scriptsize{[0.370--0.398]} & $0.142$ \scriptsize{[0.124--0.160]} & $0.545$ \\
\textbf{WGAN-GP (1D-CNN)} & $0.642$ \scriptsize{[0.630--0.654]} & $0.284$ & $0.684$ \scriptsize{[0.671--0.697]} & $0.738$ \scriptsize{[0.725--0.751]} & $0.712$ \scriptsize{[0.699--0.725]} & $0.224$ \scriptsize{[0.211--0.237]} & $0.186$ \scriptsize{[0.174--0.198]} & $0.118$ \scriptsize{[0.101--0.135]} & $0.578$ \\
\textbf{TimeGAN} & $0.689$ \scriptsize{[0.676--0.702]} & $0.378$ & $0.721$ \scriptsize{[0.708--0.734]} & $0.791$ \scriptsize{[0.778--0.804]} & $0.774$ \scriptsize{[0.761--0.787]} & $0.148$ \scriptsize{[0.137--0.159]} & $0.102$ \scriptsize{[0.091--0.113]} & $0.094$ \scriptsize{[0.078--0.110]} & $0.612$ \\
\textbf{TabDDPM-Time} & $0.618$ \scriptsize{[0.605--0.631]} & $0.236$ & $0.645$ \scriptsize{[0.632--0.658]} & \textbf{0.814} \scriptsize{[0.802--0.826]} & \textbf{0.801} \scriptsize{[0.789--0.813]} & \textbf{0.068} \scriptsize{[0.059--0.077]} & \textbf{0.049} \scriptsize{[0.041--0.057]} & \textbf{0.052} \scriptsize{[0.038--0.066]} & \textbf{0.715} \\
\bottomrule
\end{tabular}
\end{table*}

\section{Empirical Results and In-Depth Analysis}

\subsection{Global Benchmark Evaluation}
Table~\ref{tab:main_results} summarizes the global performance across all models.

\subsubsection{Key Insight 1: The Privacy-Utility Trade-off Frontier}
The empirical results reveal a steep tradeoff between downstream utility retention and membership privacy:
\begin{itemize}
    \item The \textbf{Gaussian Copula} exhibited near-ideal privacy ($\text{AUROC}_{\text{MIA}} = 0.508$), but downstream utility collapsed to $0.621$ ($95\%$ CI: $[0.608\text{--}0.634]$), severely underperforming the real-data baseline of $0.842$.
    \item \textbf{TimeGAN} captured temporal interactions effectively ($\text{TSTR} = 0.791$), but demonstrated the highest privacy vulnerability ($\text{AUROC}_{\text{MIA}} = 0.689$, adversary advantage $0.378$), overfitting to unique patient trajectories.
    \item \textbf{TabDDPM-Time} achieved the most favorable balance ($\text{TSTR} = 0.814$, $\text{AUROC}_{\text{MIA}} = 0.618$, $\Phi = 0.715$). The continuous Brownian denoising process distributes generative mass smoothly, reducing sharp memorization peaks.
\end{itemize}

\subsubsection{Key Insight 2: The Severity of Temporal Drift}
While \textbf{WGAN-GP} achieved strong static marginal alignment for Heart Rate ($W_1 = 0.042$), its transition drift was severe ($D_{\text{trans}} = 0.224$). In authentic clinical data, the probability of an instantaneous jump from Bradycardia ($<60$~bpm) to Tachycardia ($>120$~bpm) within one hour is $P_{\text{real}} = 0.003$. In WGAN-GP, this transition occurred at $P_{\text{syn}} = 0.084$---an unphysiological anomaly exceeding reality by \textbf{over 28-fold}.

Furthermore, within-stay autocorrelation decay ($D_{\text{autocorr}}$) revealed that while authentic ICU vital autocorrelation persists across 6 hours ($\rho(\tau = 6\text{h}) = 0.56$), WGAN-GP autocorrelation collapsed to near-zero within 3 hours ($\rho(\tau = 3\text{h}) = 0.05$), behaving like white noise.

\subsection{Demographic Subgroup Disparities}
Table~\ref{tab:subgroup_results} presents performance broken down across protected demographic partitions.

\begin{table}[h]
\centering
\caption{Subgroup Decomposition: Downstream Utility (TSTR AUROC) and Membership Privacy Risk (MIA AUROC).}
\label{tab:subgroup_results}
\resizebox{\columnwidth}{!}{%
\begin{tabular}{lccccc}
\toprule
\textbf{Demographic Group} & \textbf{Real (TRTR)} & \textbf{Copula} & \textbf{WGAN-GP} & \textbf{TimeGAN} & \textbf{TabDDPM} \\
\midrule
\textbf{Male} & 0.845 & 0.628 / 0.507 & 0.742 / 0.638 & 0.796 / 0.681 & 0.818 / 0.612 \\
\textbf{Female} & 0.838 & 0.612 / 0.510 & 0.724 / 0.648 & 0.782 / 0.699 & 0.809 / 0.626 \\
\textbf{Age $<$ 65} & 0.856 & 0.642 / 0.504 & 0.755 / 0.631 & 0.812 / 0.672 & 0.829 / 0.609 \\
\textbf{Age $\ge$ 65} & 0.824 & 0.589 / 0.512 & 0.698 / 0.655 & 0.748 / 0.708 & 0.792 / 0.629 \\
\textbf{White} & 0.848 & 0.631 / 0.506 & 0.749 / 0.632 & 0.802 / 0.678 & 0.821 / 0.611 \\
\textbf{Black / Afr. Am.} & 0.831 & 0.598 / 0.514 & 0.704 / 0.672 & 0.764 / 0.724 & 0.798 / 0.638 \\
\textbf{Asian} & 0.826 & 0.574 / 0.519 & 0.681 / 0.748 & 0.739 / 0.756 & 0.789 / 0.652 \\
\textbf{Hispanic / Latino} & 0.834 & 0.586 / 0.516 & 0.692 / 0.721 & 0.751 / 0.732 & 0.794 / 0.644 \\
\midrule
\textbf{Max Gap ($\Delta$)} & \textbf{0.032} & \textbf{0.068 / 0.015} & \textbf{0.074 / 0.117} & \textbf{0.074 / 0.084} & \textbf{0.040 / 0.043} \\
\bottomrule
\end{tabular}%
}
\end{table}

Under \textbf{TimeGAN} and \textbf{WGAN-GP}, demographic partitions occupying the tail of the training distribution suffered compounded harms:
\begin{enumerate}
    \item \textbf{Utility Degradation:} For Asian patients ($n = 1,642$), TSTR utility dropped under WGAN-GP to $0.681$ (compared to $0.826$ on real data), representing a penalty of over $14.5$ percentage points.
    \item \textbf{Elevated Re-identification Risk:} Concurrently, MIA vulnerability for Asian patients rose to an alarming AUROC of $0.748$ under WGAN-GP and $0.756$ under TimeGAN. Because minority patient trajectories occupy sparse regions of the clinical manifold, generative adversarial architectures overfit to these paths, effectively memorizing their trajectories.
\end{enumerate}

\section{Discussion and Clinical-Regulatory Implications}

\subsection{Regulatory Compliance: HIPAA and EU GDPR}
Our empirical findings directly impact clinical regulatory frameworks:
\begin{itemize}
    \item Under HIPAA \S164.514(b)(1) (Expert Determination Method), data release requires proving that identification risk is ``very small.'' Unconstrained generative models achieved adversary advantages up to $0.378$ ($\text{Precision@1\%FPR} = 41.2\%$), demonstrating that hundreds of training participants can be re-identified with statistical certainty.
    \item Under EU GDPR Article 9 and Recital 26, synthetic records exhibiting non-trivial membership inference advantage remain legally classified as pseudonymous personal data rather than anonymized data, exposing publishing institutions to regulatory non-compliance liability.
\end{itemize}

\subsection{Prescriptions for Generative AI Deployment}
\begin{enumerate}
    \item \textbf{Mandate Multi-Vector Stress Testing:} Reporting marginal distributions or correlation matrices alone is insufficient. Clinical evaluations must document Markovian transition drift ($D_{\text{trans}}$) and within-stay autocorrelation decay ($D_{\text{autocorr}}$).
    \item \textbf{Enforce Demographic Auditing:} Practitioners must report privacy vulnerability separately across demographic sub-cohorts to prevent disproportionate minority privacy leakage.
    \item \textbf{Adopt Score-Based Diffusion:} Continuous diffusion models (TabDDPM-Time) systematically outperform GANs across temporal realism and fairness preservation ($\Phi = 0.715$ vs. $0.578$).
\end{enumerate}

\section{Conclusion}
We introduced the \textbf{Synthetic EHR Benchmark}, an empirical framework designed to stress-test longitudinal synthetic electronic health records across Privacy, Utility, and Fairness. Our empirical investigation over 42,110 MIMIC-IV admissions reveals that unconstrained generative models trade patient privacy for predictive utility, generate biologically impossible vital sign transitions, and disproportionately expose minority cohorts to re-identification risks. We release our benchmarking library, standardized MIMIC-IV extraction routines, and evaluation protocols as open-source software to provide clinical data stewards and regulatory authorities with an adversarial gatekeeper for trustworthy generative artificial intelligence in medicine.

\section*{Declarations}
\subsection*{Ethics Approval and Consent to Participate}
This study utilized the MIMIC-IV database (v2.2), an open-access, de-identified electronic health record repository collected from Beth Israel Deaconess Medical Center. The author completed required human research training (CITI Data or Specimens Only Research, Certification ID: 183261995).

\subsection*{Data and Code Availability}
MIMIC-IV is available via PhysioNet (\url{https://physionet.org/content/mimiciv/}). The benchmarking library, cohort configs, and scripts are publicly available on GitHub at: \url{https://github.com/mtechbro94/synthetic-ehr-benchmark}.

\bibliographystyle{IEEEtran}
\bibliography{references}

\end{document}
